\documentclass[11pt]{article}

% ===================== PACKAGES =====================
\usepackage[a4paper,margin=1in]{geometry}
\usepackage{amsmath,amssymb}
\usepackage{enumitem}
\usepackage{fancyhdr}
\usepackage{xcolor} 

% ===================== HEADER & FOOTER =====================
\pagestyle{fancy}
\fancyhf{}
\lhead{CSE 400: Fundamentals of Probability in Computing}
\rhead{Lecture 6 Scribe}
\cfoot{\thepage}

% ===================== TITLE =====================
\title{
    \normalsize School of Engineering and Applied Science (SEAS), Ahmedabad University \\
    \vspace{0.2cm}
    \textbf{CSE 400: Fundamentals of Probability in Computing}\\
    \Large Lecture 6 Scribe: Discrete Random Variables \& Expectation
}
\author{Instructor: Dhaval Patel, PhD} 
\date{}

\begin{document}
\maketitle

\vspace{-1cm}
\begin{center}
    \begin{tabular}{ll}
        \textbf{Date:} & {January 22, 2025 \hspace{2.2in}} \\ [1.5ex]
    \end{tabular}
\end{center}

\hrule
\vspace{0.5cm}

% ===================== CONTENT =====================

\section{Random Variables – Motivation and Concept}
A \textbf{Random Variable (RV)} is a function that assigns a real number to each outcome in the sample space.

\subsection*{Visualization of a Random Variable}
The distribution of a random variable can be visualized as a bar diagram:
\begin{itemize}
    \item The x-axis represents the values that the random variable can take.
    \item The height of the bar at value $a_i$ is the probability $\Pr[X = a_i]$.
    \item Each probability $\Pr[X = a_i]$ is computed by evaluating the probability of the corresponding event in the sample space.
\end{itemize}

\section{Discrete and Continuous Random Variables}

\subsection{Discrete Random Variable}
A random variable is considered \textbf{discrete} if:
\begin{itemize}
    \item It has countable support.
    \item It is described using a \textbf{Probability Mass Function (PMF)}.
    \item Probabilities are assigned to single values.
    \item Each possible value has strictly positive probability.
\end{itemize}

\subsection{Continuous Random Variable}
A random variable is considered \textbf{continuous} if:
\begin{itemize}
    \item It has uncountable support.
    \item It is described using a \textbf{Probability Density Function (PDF)}.
    \item Probabilities are assigned to \textit{intervals} of values.
    \item Each individual value has zero probability.
\end{itemize}

\section{Example 1 – Tossing 3 Fair Coins}
\textbf{Problem Setup:} The experiment consists of tossing 3 fair coins. Let $Y$ denote the number of heads that appear.
The possible values of $Y$ are $Y \in \{0, 1, 2, 3\}$.

\subsection*{Computation of Probabilities}
Each single outcome of 3 fair coin tosses (e.g., $h,t,h$) has a probability of $\frac{1}{8}$.

\begin{itemize}
    \item \textbf{P(Y = 0):} $\{t,t,t\}$
    \[ P(Y = 0) = \frac{1}{8} \]
    
    \item \textbf{P(Y = 1):} $\{t,t,h\}, \{t,h,t\}, \{h,t,t\}$
    \[ P(Y = 1) = \frac{3}{8} \]
    
    \item \textbf{P(Y = 2):} $\{t,h,h\}, \{h,t,h\}, \{h,h,t\}$
    \[ P(Y = 2) = \frac{3}{8} \]
    
    \item \textbf{P(Y = 3):} $\{h,h,h\}$
    \[ P(Y = 3) = \frac{1}{8} \]
\end{itemize}

\subsection*{Verification of Total Probability}
Since $Y$ must take one of the values $\{0,1,2,3\}$:
\[ \sum_{i=0}^{3} P(Y=i) = \frac{1}{8} + \frac{3}{8} + \frac{3}{8} + \frac{1}{8} = 1 \]

\section{Probability Mass Function (PMF)}
\textbf{Definition:} A random variable $X$ is discrete if it takes on at most a countable number of possible values.
Range: $R_X = \{x_1, x_2, x_3, \dots\}$ (finite or countably infinite).

The function $P_X(x_k) = P(X = x_k)$ for $k = 1,2,3,\dots$ is called the \textbf{Probability Mass Function (PMF)} of $X$.

\textbf{Property:}
\[ \sum_{k=1}^{\infty} P_X(x_k) = 1 \]

\section{PMF – Worked Example}
\textbf{Problem:} The PMF of a random variable $X$ is given by:
\[ p(i) = c \frac{\lambda^i}{i!}, \quad i = 0,1,2,\dots \quad (\text{where } \lambda > 0) \]
Find $P(X=0)$ and $P(X > 2)$.

\subsection*{Step 1: Determine the Constant $c$}
Using the property that total probability sums to 1:
\[ \sum_{i=0}^{\infty} p(i) = 1 \implies \sum_{i=0}^{\infty} c \frac{\lambda^i}{i!} = 1 \]
\[ c \sum_{i=0}^{\infty} \frac{\lambda^i}{i!} = 1 \]
Using the Taylor series expansion $e^{\lambda} = \sum_{i=0}^{\infty} \frac{\lambda^i}{i!}$:
\[ c e^{\lambda} = 1 \implies c = e^{-\lambda} \]

\subsection*{Step 2: Compute $P(X = 0)$}
\[ P(X = 0) = c \frac{\lambda^0}{0!} = e^{-\lambda} \cdot 1 = e^{-\lambda} \]

\subsection*{Step 3: Compute $P(X > 2)$}
\[ P(X > 2) = 1 - P(X \le 2) \]
\[ P(X \le 2) = P(X=0) + P(X=1) + P(X=2) \]
\[ P(X \le 2) = e^{-\lambda} + e^{-\lambda}\frac{\lambda^1}{1!} + e^{-\lambda}\frac{\lambda^2}{2!} \]
\[ P(X > 2) = 1 - e^{-\lambda} \left( 1 + \lambda + \frac{\lambda^2}{2} \right) \]

\section{Bayes’ Theorem – Recap}
Using the definition of conditional probability $\Pr(A \cap B_i) = \Pr(B_i \mid A)\Pr(A)$, we obtain Bayes' Formula:
\[ \Pr(B_i \mid A) = \frac{\Pr(A \mid B_i)\Pr(B_i)}{\sum_{j=1}^{n} \Pr(A \mid B_j)\Pr(B_j)} \]
\begin{itemize}
    \item $\Pr(B_i)$: \textbf{A priori} probability
    \item $\Pr(B_i \mid A)$: \textbf{Posterior} probability
\end{itemize}

\section{Bayes’ Theorem – Example 1: Auditorium}
\textbf{Problem:} An auditorium has 30 rows.
\begin{itemize}
    \item Row $k$ has $k+10$ seats (Row 1 has 11, Row 30 has 40).
    \item A door prize is given by randomly selecting a row (equally likely, $1/30$), then a seat within that row.
\end{itemize}
Let $S_{15}$ be the event that Seat 15 is chosen. Let $R_k$ be the event Row $k$ is chosen.
Note: Seat 15 only exists in rows where total seats $\ge 15$. Since seats $= k+10$, we need $k+10 \ge 15 \implies k \ge 5$.

\subsection*{1. Computation of $P(S_{15})$}
\[ P(S_{15}) = \sum_{k=5}^{30} P(S_{15} \mid R_k) P(R_k) \]
\[ = \sum_{k=5}^{30} \left( \frac{1}{k+10} \right) \cdot \frac{1}{30} \]
\[ = \frac{1}{30} \sum_{k=5}^{30} \frac{1}{k+10} \approx 0.0342 \]

\subsection*{2. Conditional Probability $P(R_{20} \mid S_{15})$}
Given that Seat 15 was picked, what is the probability it came from Row 20?
\[ P(R_{20} \mid S_{15}) = \frac{P(S_{15} \mid R_{20})P(R_{20})}{P(S_{15})} \]
\[ = \frac{\left(\frac{1}{30}\right) \cdot \left(\frac{1}{30}\right)}{0.0342} \approx 0.0325 \]

\section{Bayes’ Theorem – Example 2: Communication System}
A communication system sends binary data (0 or 1). The receiver $RX$ may make mistakes (0 detected as 1, or 1 detected as 0). The system is modeled using conditional probabilities of detection given the transmitted bit.

\bigskip
\hrule
\vspace{0.2cm}
\begin{center}
    \small \textit{End of Lecture 6 Scribe}
\end{center}

\end{document}